\documentclass[10pt,twocolumn]{article}

\usepackage[margin=0.78in,top=0.62in,bottom=0.70in]{geometry}
\usepackage[T1]{fontenc}
\usepackage[utf8]{inputenc}
\usepackage{lmodern}
\usepackage{amsmath,amssymb,mathtools}
\usepackage{booktabs,array,multirow}
\usepackage{enumitem}
\usepackage{xcolor}
\usepackage{microtype}
\usepackage[hidelinks]{hyperref}
\usepackage{titlesec}

\setlength{\columnsep}{22pt}
\setlength{\parindent}{0pt}
\setlength{\parskip}{2.2pt}
\pagestyle{plain}

\makeatletter
\renewcommand\normalsize{\@setfontsize\normalsize{9pt}{11.2pt}}
\normalsize
\makeatother

\titleformat{\section}{\large\bfseries}{}{0em}{}[\vspace{0.18em}\titlerule]
\titleformat{\subsection}{\normalsize\bfseries}{}{0em}{}
\titleformat{\subsubsection}{\normalsize\itshape}{}{0em}{}
\titlespacing*{\section}{0pt}{1.2ex}{0.75ex}
\titlespacing*{\subsection}{0pt}{0.8ex}{0.35ex}
\titlespacing*{\subsubsection}{0pt}{0.45ex}{0.25ex}

\setlist[itemize]{leftmargin=1.15em,itemsep=0.3ex,topsep=0.35ex,parsep=0pt,partopsep=0pt}
\setlist[enumerate]{leftmargin=1.35em,itemsep=0.3ex,topsep=0.35ex,parsep=0pt,partopsep=0pt}
\setlist[description]{leftmargin=1.15em,itemsep=0.3ex,topsep=0.35ex,parsep=0pt,partopsep=0pt,font=\normalfont\itshape}

\newcommand{\bbR}{\mathbb{R}}
\newcommand{\x}{\mathbf{x}}
\newcommand{\y}{\mathbf{y}}
\newcommand{\w}{\mathbf{w}}
\newcommand{\E}{\mathcal{E}}
\newcommand{\Hc}{\mathcal{H}}
\newcommand{\argmin}{\mathop{\mathrm{arg\,min}}}
\newcommand{\argmax}{\mathop{\mathrm{arg\,max}}}
\newcommand{\defeq}{\vcentcolon=}
\newcommand{\chead}[1]{\vspace{0.15em}\textbf{#1:} }
\newcommand{\tightdisplay}{\vspace{-0.5em}}

\begin{document}
\thispagestyle{plain}
\begin{center}
{\LARGE\bfseries 15-288 Machine Learning Cheatsheet}\\[-0.15em]
{\small Maximally condensed summary of the notes}
\end{center}
\vspace{-0.8em}

\section{Core Setup}
\chead{Learning problem} From samples $\{(\x^{(i)},y^{(i)})\}_{i=1}^N$, learn $f:X\to Y$ that minimizes error on unseen data.\

\chead{Tasks}
\begin{itemize}
  \item \textbf{Classification:} $Y$ discrete; predict a class.
  \item \textbf{Regression:} $Y\subseteq\bbR^m$ continuous; predict a value.
  \item \textbf{Unsupervised:} no labels; cluster, compress, discover structure.
  \item \textbf{RL:} learn a policy maximizing cumulative reward.
\end{itemize}
\chead{Pipeline} Define task/metric $\to$ collect and clean data $\to$ EDA $\to$ engineer features $\to$ choose hypothesis class $\to$ train $\to$ validate/tune $\to$ test once $\to$ deploy/monitor.

\subsection{Notation and Generalization}
\begin{itemize}
  \item \textbf{Hypothesis class} $\Hc$: candidate functions; its shape is the model bias.
  \item \textbf{Loss} $\ell(\hat y,y)$: error on one example.
  \item \textbf{Empirical risk} $L_{\mathrm{emp}}(f)=\frac{1}{N}\sum_{i=1}^N \ell(f(\x^{(i)}),y^{(i)})$.
  \item \textbf{Generalization} means low expected loss on unseen samples from the same distribution.
  \item \textbf{Inductive bias} is necessary: many hypotheses fit training data; we prefer simpler/useful ones.
\end{itemize}
\chead{Train/val/test} Fit parameters on train, tune choices on validation/CV, report final performance once on test. Never let test influence training.

\subsection{Paradigm Comparison}
\begin{center}
\renewcommand{\arraystretch}{1.12}
\begin{tabular}{>{\raggedright\arraybackslash}p{0.19\linewidth} >{\raggedright\arraybackslash}p{0.24\linewidth} >{\raggedright\arraybackslash}p{0.24\linewidth} >{\raggedright\arraybackslash}p{0.23\linewidth}}
\toprule
 & Supervised & Unsupervised & Reinforcement \\
\midrule
Data & labeled $(\x,y)$ & unlabeled $\x$ only & generated by interaction \\
Signal & label error & similarity/structure & reward/cost \\
Goal & predict $y$ & discover structure & learn policy \\
Examples & class/regression & clustering, PCA, density & control, games, robotics \\
\bottomrule
\end{tabular}
\end{center}
\chead{Label scarcity} Labels are expensive, slow, noisy, and often require experts; that motivates unsupervised, self-supervised, and transfer-learning pipelines.

\section{Supervised Learning Essentials}
\subsection{Losses and Metrics}
\chead{Classification losses}
\begin{itemize}
  \item 0/1 loss: $\ell(\hat y,y)=\mathbf{1}[\hat y\neq y]$.
  \item Cross-entropy (binary): $-\big(y\log p + (1-y)\log(1-p)\big)$.
  \item Cross-entropy (multi-class): $-\sum_k y_k \log p_k$.
  \item Hinge loss: $\max(0,1-y(\w^\top\x+b))$ for $y\in\{-1,+1\}$.
\end{itemize}

\newpage

\chead{Regression losses}
\begin{itemize}
  \item SSE $=\sum_i (y_i-\hat y_i)^2$.
  \item MSE $=\frac{1}{N}\sum_i (y_i-\hat y_i)^2$.
  \item RMSE $=\sqrt{\text{MSE}}$.
  \item MAE $=\frac{1}{N}\sum_i |y_i-\hat y_i|$; more robust to outliers.
\end{itemize}

\subsection{Confusion-Matrix Metrics}
For binary classification:
\[
\text{Accuracy}=\frac{TP+TN}{TP+TN+FP+FN},\qquad
\text{Precision}=\frac{TP}{TP+FP}
\]
\[
\text{Recall}=\frac{TP}{TP+FN},\qquad
F_1=\frac{2PR}{P+R}.
\]
Use \textbf{accuracy} for balanced classes, \textbf{precision} when false positives hurt, \textbf{recall} when false negatives hurt, \textbf{$F_1$} for imbalanced positive-class detection, \textbf{ROC-AUC/PR-AUC} when threshold-free ranking matters.

\subsection{Overfitting and Model Selection}
\begin{itemize}
  \item \textbf{Overfitting:} training error low but validation/test error high.
  \item \textbf{Underfitting:} both errors high; model too simple or badly trained.
  \item \textbf{Bias--variance tradeoff:} more flexible models reduce bias but raise variance.
  \item Choose complexity at the validation minimum, not at best training score.
\end{itemize}
\chead{Generalization gap} $L_{\mathrm{gen}}-L_{\mathrm{emp}}$ is the quantity we care about, but true $L_{\mathrm{gen}}$ is unknown; validation/test error estimates it.
\chead{Promote generalization}
\begin{itemize}
  \item Match model capacity to the data size/noise level.
  \item Add regularization terms $\lambda R(\w)$ or early stopping.
  \item Gather more representative data when possible.
  \item Use correct train/val/test separation and robust validation protocols.
\end{itemize}

\subsection{Margin-Based View of Classification}
For binary labels $y\in\{-1,+1\}$ and score $f(\x)$, the margin is
\[
\mu(\x)=y\,f(\x).
\]
\begin{itemize}
  \item $\mu>0$: correct sign; $\mu<0$: misclassified.
  \item Larger positive margin means more confident correct prediction.
  \item 0/1, hinge, logistic, and exponential losses are all surrogate losses on $\mu$.
\end{itemize}

\section{Validation and Workflow}
\subsection{Hold-Out and Cross-Validation}
\begin{itemize}
  \item \textbf{Hold-out:} split once (e.g. train/val/test); fast but noisy.
  \item \textbf{$K$-fold CV:} split into $K$ folds, train on $K-1$, validate on the last, average scores.
  \item Typical choice: $K=5$ or $10$.
  \item Use CV for model/hyperparameter selection, then retrain on full train set.
\end{itemize}
\chead{Leakage rule} Any operation that learns from data (imputation, scaling, PCA, feature selection, target encoding) must be fit on training folds only, then applied to validation/test.

\subsection{Design Checklist}
\begin{enumerate}
  \item Define output space, metric, and operating constraints first.
  \item Pick a hypothesis class with the right inductive bias.
  \item Choose preprocessing compatible with the model.
  \item Select a loss aligned with the business/scientific objective.
  \item Tune hyperparameters only with validation/CV.
  \item Interpret results through errors, calibration, and failure cases.
\end{enumerate}

\subsection{Hypothesis Classes}
\begin{itemize}
  \item \textbf{Parametric:} finite parameters, fixed form (linear/logistic regression, neural nets with fixed architecture).
  \item \textbf{Nonparametric:} complexity grows with data ($k$-NN, trees).
  \item \textbf{Consistency} on train does \emph{not} imply generalization.
\end{itemize}

\section{Data Preprocessing}
\subsection{Outliers}
\chead{Why care} Outliers distort means, variances, correlations, distances, regression slopes, and gradient updates.
\begin{itemize}
  \item \textbf{Z-score:} $z=\frac{x-\mu}{\sigma}$; flag if $|z|>3$ (assumes near-Gaussian data).
  \item \textbf{Modified Z-score:} $0.6745\frac{x-\text{median}}{\text{MAD}}$; robust.
  \item \textbf{IQR rule:} IQR $=Q_3-Q_1$; outlier if $x<Q_1-1.5\,\text{IQR}$ or $x>Q_3+1.5\,\text{IQR}$.
  \item \textbf{Percentile clipping/winsorization:} cap extremes at chosen quantiles.
\end{itemize}
\chead{Handle by} fixing data-entry errors, deleting impossible points, clipping, transforming, robust scaling, or using robust models. Remove only with reason.

\subsection{Scaling}
\[
\text{Standardize: }x'=\frac{x-\mu}{\sigma}
\qquad
\text{Min--max: }x'=\frac{x-x_{\min}}{x_{\max}-x_{\min}}
\]
\[
\text{Robust scale: }x'=\frac{x-\text{median}}{Q_3-Q_1}
\qquad
\text{MaxAbs: }x'=\frac{x}{\max|x|}
\]
\chead{Needed for} distance-based methods ($k$-NN, clustering), margin/dot-product methods (SVM, logistic regression), gradient-based models, PCA, neural nets. Usually not needed for trees/random forests.

\subsection{Feature Types and Encoding}
\begin{itemize}
  \item \textbf{Numerical:} continuous/discrete quantities.
  \item \textbf{Categorical nominal:} no order; use one-hot.
  \item \textbf{Categorical ordinal:} order exists; integer/ordinal encoding can be valid.
  \item \textbf{Binary:} 0/1 directly.
\end{itemize}
\chead{Label encoding} Maps categories to integers; dangerous for nominal categories because it injects fake ordering.\

\chead{One-hot encoding} Adds one binary feature per category. Avoid exact linear dependence by dropping one level in strict linear models (dummy-variable trap).\

\chead{Target encoding} Replace category by target mean; powerful but leakage-prone, so compute only within training folds.

\subsection{Correlation and Feature Engineering}
Pearson correlation:
\[
r_{xy}=\frac{\sum_i (x_i-\bar x)(y_i-\bar y)}{\sqrt{\sum_i (x_i-\bar x)^2}\sqrt{\sum_i (y_i-\bar y)^2}}\in[-1,1].
\]
\begin{itemize}
  \item $r\approx 1$: strong positive linear relation; $r\approx -1$: strong negative; $r\approx 0$: no linear relation.
  \item Correlation \emph{is not} causation and misses nonlinear dependence.
  \item Highly correlated features can create multicollinearity in linear models.
  \item Polynomial features allow linear models to fit nonlinear trends by expanding inputs (e.g. $x,x^2,x^3,xy$).
\end{itemize}
\chead{Practical rule} Use correlation matrices to remove redundant numeric features, but do not drop features solely by low correlation if nonlinear models or interactions may matter.

\section{Image Data and Classical Vision Features}
\subsection{Image Representation}
An image is a tensor: grayscale $H\times W$, color $H\times W\times C$ with channels such as RGB. Pixel intensities are features; resolution, channel order, and dtype matter.
\begin{itemize}
  \item \textbf{Metadata (EXIF):} timestamps, camera settings, orientation, GPS; can be informative or leakage-prone.
  \item \textbf{Color spaces:} RGB for display, grayscale for intensity, HSV for hue/saturation separation.
  \item \textbf{Global statistics:} channel means, variances, min/max, entropy.
\end{itemize}

\subsection{Histograms and Patches}
\begin{itemize}
  \item \textbf{Color histograms:} count pixel intensities per bin; robust to small translations but lose exact spatial layout.
  \item \textbf{Spatial patching:} compute local features on tiles to preserve rough spatial structure.
  \item \textbf{Combined descriptors:} concatenate global color/texture/edge summaries into one vector.
\end{itemize}

\subsection{Filtering and Convolution}
For image $I$ and kernel $K$, cross-correlation at location $(i,j)$ is the weighted sum of a local patch. Convolution flips the kernel before sliding; in deep learning, people often say ``convolution'' while computing cross-correlation.
\begin{itemize}
  \item \textbf{Padding} preserves border size; \textbf{stride} controls step size.
  \item Common kernels: blur/smoothing, sharpen, edge detectors, emboss.
  \item Linear filters combine by addition; median filters are nonlinear and robust to salt-and-pepper noise.
\end{itemize}

\subsection{Edges, Gradients, and HOG}
Image gradients measure intensity change:
\[
\nabla I = \left(\frac{\partial I}{\partial x},\frac{\partial I}{\partial y}\right),
\qquad
\|\nabla I\|=\sqrt{I_x^2+I_y^2},
\qquad
\theta=\tan^{-1}(I_y/I_x).
\]
\begin{itemize}
  \item \textbf{Sobel operators} approximate horizontal/vertical derivatives.
  \item \textbf{Canny pipeline:} smooth $\to$ gradient $\to$ non-max suppression $\to$ hysteresis thresholding.
  \item \textbf{HOG:} bin local edge orientations, weight by magnitude, normalize over blocks.
  \item HOG is useful for classical object detection and shape-sensitive tasks.
\end{itemize}

\subsection{Thresholding, Corners, and Matching}
\begin{itemize}
  \item \textbf{Global thresholding:} one cutoff for all pixels.
  \item \textbf{Adaptive thresholding:} threshold depends on local neighborhood.
  \item \textbf{Otsu:} choose threshold maximizing between-class variance.
  \item \textbf{Corners:} points with strong intensity variation in two directions (e.g. Harris detector).
  \item \textbf{SIFT/ORB:} keypoint descriptors for matching across scale/rotation/view changes.
  \item \textbf{Template matching:} slide a patch and compute similarity; simple but sensitive to scale/pose/illumination.
\end{itemize}
\chead{Classical image pipeline} preprocess $\to$ extract descriptors (histogram/HOG/keypoints) $\to$ normalize $\to$ feed into SVM/logistic/tree model.

\section{Linear Regression and OLS}
\subsection{Model and Matrix Form}
Linear model:
\[
\hat y = w_0 + \sum_{j=1}^d w_j x_j = \w^\top \tilde{\x},
\]
where $\tilde{\x}$ includes a leading 1 for the bias term. With design matrix $X\in\bbR^{N\times (d+1)}$,
\[
\hat{\y}=X\w,\qquad L(\w)=\|X\w-\y\|_2^2.
\]
\chead{Interpretation} Linear in parameters, not necessarily in raw inputs (polynomial features keep model linear in $\w$).

\subsection{Normal Equations}
For least squares,
\[
\nabla_{\w}L = 2X^\top(X\w-\y)=0
\quad\Rightarrow\quad
X^\top X\w = X^\top\y.
\]
If $X^\top X$ is invertible,
\[
\hat{\w}=(X^\top X)^{-1}X^\top\y.
\]
\begin{itemize}
  \item MSE is convex in $\w$ for linear regression.
  \item Closed form is elegant but costly/unstable for very large or collinear features.
  \item Gradient methods scale better in modern settings.
\end{itemize}

\subsection{Regularization}
\begin{itemize}
  \item \textbf{Ridge:} minimize $\|X\w-\y\|_2^2 + \lambda\|\w\|_2^2$; shrinks weights smoothly.
  \item \textbf{Lasso:} minimize $\|X\w-\y\|_2^2 + \lambda\|\w\|_1$; can force exact zeros.
  \item \textbf{Elastic Net:} combine $L_1$ and $L_2$.
\end{itemize}
Regularization reduces variance, combats multicollinearity, and improves generalization.

\subsection{Linear Models as Feature Machines}
\begin{itemize}
  \item Linear decision functions can still model nonlinear boundaries after a feature map $\phi(\x)$.
  \item Basis expansions: polynomials, radial basis functions, handcrafted image descriptors, embeddings.
  \item Most of classical ML can be seen as: engineer a representation first, then fit a simpler model on top.
\end{itemize}

\section{$k$-Nearest Neighbors}
\chead{Idea} Store training set; for a query point, find the $k$ closest examples.
\begin{itemize}
  \item \textbf{Classification:} majority vote among neighbors.
  \item \textbf{Regression:} average/weighted average neighbor targets.
  \item Common distances: Euclidean $\|\x-\x'\|_2$, Manhattan $\|\x-\x'\|_1$.
\end{itemize}
\chead{Hyperparameter} Small $k$ $\Rightarrow$ low bias/high variance; large $k$ $\Rightarrow$ smoother boundary, higher bias/lower variance.
\begin{itemize}
  \item Requires scaling.
  \item Prediction is expensive; training is trivial.
  \item Suffers in high dimension (distance concentration / curse of dimensionality).
\end{itemize}
\chead{Weighted neighbors} Distance-weighted voting/regression gives nearer points more influence and often smooths ties.

\section{Optimization and Gradient Methods}
\subsection{Gradient Descent}
Update rule:
\[
\w_{t+1}=\w_t-\eta\nabla_{\w}L(\w_t),
\]
where $\eta$ is the learning rate.
\begin{itemize}
  \item Too small $\eta$: slow convergence.
  \item Too large $\eta$: oscillation/divergence.
  \item Stop by validation plateau, gradient norm, or max epochs.
\end{itemize}
\chead{Variants}
\begin{itemize}
  \item \textbf{Batch GD:} one gradient over full dataset; stable but slow.
  \item \textbf{SGD:} one example at a time; noisy but cheap and can escape shallow regions.
  \item \textbf{Mini-batch GD:} practical default; vectorized and reasonably stable.
\end{itemize}

\subsection{Advanced Optimizers}
\begin{itemize}
  \item \textbf{Momentum:} accumulates velocity to smooth descent.
  \item \textbf{RMSProp:} adapts step size using running squared gradients.
  \item \textbf{Adam/AdamW:} momentum + adaptive scaling; strong default for deep learning.
\end{itemize}
\chead{Deep-learning defaults} mini-batches, shuffled data, Adam/AdamW, validation monitoring, early stopping, learning-rate schedules.

\subsection{Optimization Pitfalls}
\begin{itemize}
  \item Nonconvex objectives can contain saddles and many local minima.
  \item Ill-conditioned curvature causes zig-zagging and slow training.
  \item Feature scaling and normalization often improve optimization as much as model choice.
\end{itemize}

\section{Linear Classification}
\subsection{Linear Decision Functions}
General linear classifier:
\[
g(\x)=\w^\top\x+b,
\qquad
\hat y = \mathrm{sign}(g(\x)) \text{ or } \argmax_k g_k(\x).
\]
Decision boundary is the hyperplane $\w^\top\x+b=0$. Linear models work well when classes are approximately linearly separable or after good feature engineering.

\subsection{Logistic Regression}
Binary logistic regression models
\[
p(y=1\mid\x)=\sigma(z),\qquad z=\w^\top\x+b,\qquad \sigma(z)=\frac{1}{1+e^{-z}}.
\]
\[
\log\frac{p}{1-p}=\w^\top\x+b
\]
is the log-odds (logit). Predict class 1 when $p\ge \tau$ (often $\tau=0.5$).

\chead{Loss} Minimize binary cross-entropy, often plus regularization.\

\chead{Why likelihood} Logistic regression is equivalent to maximum conditional likelihood for Bernoulli labels.

\chead{Interpretation} Increasing feature $x_j$ by 1 changes log-odds by $w_j$ (holding others fixed).\

\chead{Thresholding} Decision threshold need not be $0.5$; raise threshold for higher precision, lower it for higher recall.

\chead{Multi-class} Softmax regression:
\[
p(y=k\mid\x)=\frac{e^{z_k}}{\sum_j e^{z_j}}, \qquad z_k=\w_k^\top \x+b_k.
\]
\[
\ell_{\mathrm{CE}}=-\sum_k y_k\log p_k.
\]

\subsection{SVM}
\chead{Geometric margin} For hyperplane $(\w,b)$ and example $(\x_i,y_i)$ with $y_i\in\{-1,+1\}$:
\[
\gamma_i = \frac{y_i(\w^\top\x_i+b)}{\|\w\|}.
\]
Support vectors are points lying on/inside the margin and determine the boundary.

\chead{Hard-margin SVM}
\[
\min_{\w,b}\;\frac12\|\w\|^2 \quad \text{s.t.}\quad y_i(\w^\top\x_i+b)\ge 1\ \forall i.
\]

\chead{Soft-margin SVM}
\[
\min_{\w,b,\xi}\;\frac12\|\w\|^2 + C\sum_i \xi_i
\quad \text{s.t.}\quad
y_i(\w^\top\x_i+b)\ge 1-\xi_i,\; \xi_i\ge 0.
\]
Equivalent viewpoint: regularized hinge-loss minimization.

\chead{Kernel trick} Replace inner products by $K(\x,\x')=\phi(\x)^\top\phi(\x')$ to get nonlinear boundaries without explicitly constructing $\phi$. Common kernels: linear, polynomial, RBF.\

\chead{When kernels help} Small/medium datasets with nonlinear boundaries and meaningful similarity structures; usually less attractive at very large scale.

\chead{Logistic vs. SVM} Logistic gives calibrated probabilities; SVM optimizes margin and often sharper boundaries.

\section{Decision Trees}
\chead{Idea} Recursively partition feature space into axis-aligned regions, predict by leaf majority class or leaf mean.

\subsection{Split Criteria}
\begin{itemize}
  \item \textbf{Entropy:} $H(S)=-\sum_c p_c\log_2 p_c$.
  \item \textbf{Information gain:} $IG(S,\text{split}) = H(S)-\sum_j \frac{|S_j|}{|S|}H(S_j)$.
  \item \textbf{Gini impurity:} $G(S)=1-\sum_c p_c^2$.
  \item \textbf{Regression trees:} choose split with largest variance/MSE reduction.
\end{itemize}
\chead{Continuous features} Test thresholds of the form $x_j\le t$ vs. $x_j>t$.

\subsection{Pros, Cons, and Pruning}
\begin{itemize}
  \item Pros: interpretable, handles mixed features, no scaling required, nonlinear boundaries.
  \item Cons: high variance, greedy, unstable to data perturbations, can overfit deeply.
  \item \textbf{Pre-pruning:} max depth, min samples split/leaf, max leaves.
  \item \textbf{Post-pruning:} cost-complexity pruning trims weak subtrees.
\end{itemize}
\chead{Ensembles} Bagging/random forests reduce variance via averaging many trees; boosting fits sequentially to residual errors.
\chead{Feature importance caveat} Tree importances can be biased toward high-cardinality features; permutation importance is often safer.

\section{Unsupervised Learning}
\subsection{Clustering}
\chead{Goal} Partition unlabeled data into groups so points in the same cluster are more similar than points in different clusters.
\begin{itemize}
  \item Choice of similarity metric is an inductive bias.
  \item Number of clusters is often unknown and must be selected/estimated.
  \item Cluster labels are arbitrary; cluster quality depends on the problem definition.
\end{itemize}
\chead{$k$-means} Alternate between assigning each point to nearest centroid and recomputing centroids as cluster means; minimizes within-cluster SSE.
\begin{itemize}
  \item Fast and common, but assumes roughly spherical clusters under Euclidean distance.
  \item Sensitive to scaling, initialization, and outliers.
\end{itemize}

\subsection{Dimensionality Reduction}
\begin{itemize}
  \item High-dimensional data often lies near a lower-dimensional manifold.
  \item \textbf{PCA:} find orthogonal directions of maximum variance; useful for compression, denoising, visualization, de-correlation.
  \item Must scale features appropriately before PCA when units differ.
\end{itemize}

\subsection{Generative Modeling View}
\begin{itemize}
  \item Generative models learn/discuss the data distribution $p(\x)$ or joint structure in data.
  \item Useful for synthesis, density estimation, compression, anomaly detection, and representation learning.
\end{itemize}

\section{Neural Networks}
\subsection{Perceptrons and MLPs}
One neuron computes
\[
a = \w^\top\x+b,\qquad h=\phi(a),
\]
where $\phi$ is an activation function. Stacking layers gives an MLP:
\[
\x \to h^{(1)} \to h^{(2)} \to \cdots \to \hat y.
\]
\begin{itemize}
  \item Without nonlinear activations, stacked linear layers collapse to one linear map.
  \item Hidden layers learn intermediate features.
  \item Universal approximation: sufficiently wide networks can approximate many functions, but optimization/generalization still matter.
\end{itemize}

\subsection{Activations and Output Layers}
\begin{itemize}
  \item \textbf{Sigmoid:} $(1+e^{-z})^{-1}$; bounded, but saturates and causes vanishing gradients.
  \item \textbf{Tanh:} zero-centered but still saturates.
  \item \textbf{ReLU:} $\max(0,z)$; simple and dominant default for hidden layers.
  \item \textbf{Leaky ReLU/GELU:} help with dead neurons/smoother behavior.
  \item \textbf{Output:} identity for regression, sigmoid for binary class, softmax for multi-class.
\end{itemize}

\subsection{Backpropagation}
Backprop computes gradients layer-by-layer using the chain rule. For layer $\ell$:
\[
\delta^{(\ell)} = \frac{\partial L}{\partial a^{(\ell)}}
= \left(W^{(\ell+1)\top}\delta^{(\ell+1)}\right)\odot \phi'(a^{(\ell)}).
\]
Then
\[
\frac{\partial L}{\partial W^{(\ell)}} = \delta^{(\ell)} h^{(\ell-1)\top},
\qquad
\frac{\partial L}{\partial b^{(\ell)}} = \delta^{(\ell)}.
\]

\subsection{Training Dynamics}
\begin{itemize}
  \item \textbf{Epoch:} one full pass through training data.
  \item \textbf{Batch size:} examples processed before one update.
  \item Problems: vanishing/exploding gradients, overfitting, poor conditioning.
  \item Fixes: normalization, good initialization, ReLU-like activations, residual ideas, regularization.
\end{itemize}
\chead{Regularization in NNs} weight decay, dropout, data augmentation, early stopping, batch/layer normalization.

\section{Convolutional Neural Networks}
\subsection{Why CNNs}
MLPs ignore image locality and use too many parameters on pixels. CNNs exploit:
\begin{itemize}
  \item \textbf{Local connectivity:} filters inspect small receptive fields.
  \item \textbf{Weight sharing:} one filter scans all spatial positions.
  \item \textbf{Translation-sensitive feature maps:} same detector reused across image.
\end{itemize}

\subsection{Main Layers}
\begin{itemize}
  \item \textbf{Convolution:} sliding kernels produce feature maps.
  \item \textbf{Activation:} usually ReLU/GELU.
  \item \textbf{Pooling:} downsample (max/avg) to reduce spatial size and add invariance.
  \item \textbf{Flatten / global pooling + dense head:} final classification/regression stage.
\end{itemize}
\chead{Parameter efficiency} A small set of filters can replace enormous fully connected layers by reusing weights over space.

\subsection{Transfer Learning}
Use a pretrained CNN as:
\begin{itemize}
  \item \textbf{Fixed feature extractor:} freeze backbone, train new head.
  \item \textbf{Fine-tuned model:} unfreeze some/all deeper layers and adapt to new data.
\end{itemize}
Especially useful when target dataset is small.

\section{Modern NLP and Generative AI}
\subsection{Basic NLP Representations}
\begin{itemize}
  \item \textbf{Tokenization:} split text into tokens (character, word, or subword).
  \item \textbf{Subword tokenization} is standard: balances vocabulary size and OOV handling.
  \item \textbf{BoW/one-hot:} sparse, ignore order.
  \item \textbf{Embeddings:} dense vectors that encode semantic similarity.
\end{itemize}

\subsection{Transformer Core}
Transformers replace recurrence with attention and parallel sequence processing.
\chead{Scaled dot-product attention}
\[
\mathrm{Attention}(Q,K,V)=\mathrm{softmax}\!\left(\frac{QK^\top}{\sqrt{d_k}}\right)V.
\]
\begin{itemize}
  \item \textbf{Queries, keys, values} are learned projections of token embeddings.
  \item \textbf{Positional encoding} injects order information.
  \item \textbf{Residual + normalization} stabilize training.
  \item \textbf{Feed-forward blocks} add nonlinear transformation per token.
\end{itemize}

\subsection{Transformer Families}
\begin{itemize}
  \item \textbf{Encoder-only (BERT):} bidirectional context; best for understanding/classification.
  \item \textbf{Decoder-only (GPT):} causal masking; predicts next token autoregressively.
  \item \textbf{Encoder--decoder (T5):} seq2seq tasks like translation/summarization.
\end{itemize}
\chead{Typical workflow} tokenize $\to$ map to embeddings $\to$ run model $\to$ fine-tune/use pretrained head $\to$ evaluate with task metric.

\subsection{Generative Models}
\begin{itemize}
  \item \textbf{Autoencoder:} encoder compresses to latent code, decoder reconstructs; optimize reconstruction loss.
  \item \textbf{GAN:} generator tries to fool discriminator; discriminator distinguishes real vs fake.
  \item \textbf{LLM generation controls:} temperature controls randomness; top-$k$/top-$p$ restrict candidate next tokens.
\end{itemize}

\section{High-Yield Comparisons}
\subsection{Algorithm Selection Heuristics}
\begin{center}
\renewcommand{\arraystretch}{1.15}
\begin{tabular}{>{\raggedright\arraybackslash}p{0.21\linewidth} >{\raggedright\arraybackslash}p{0.34\linewidth} >{\raggedright\arraybackslash}p{0.34\linewidth}}
\toprule
Model & Strengths & Weaknesses \\
\midrule
Linear/Logistic & fast, interpretable, strong baseline & misses complex nonlinear patterns \\
$k$-NN & simple, nonparametric, flexible boundary & slow prediction, scaling sensitive, weak in high $d$ \\
SVM & strong margins, effective in medium-sized/high-$d$ data, kernels & expensive on large $N$, less interpretable \\
Trees & interpretable, mixed data, no scaling & unstable, overfit alone \\
Ensembles & strong accuracy, robust & less interpretable, heavier compute \\
MLP & flexible universal approximator & tuning/data hungry, weak inductive bias on images/text \\
CNN & excellent on images/spatial data & architecture/training complexity \\
Transformers & dominant on sequence/text, scalable & memory/compute heavy \\
\bottomrule
\end{tabular}
\end{center}

\subsection{Common Failure Modes}
\begin{itemize}
  \item \textbf{Leakage:} preprocessing fitted on all data, peeking at test labels, target leakage from future information.
  \item \textbf{Metric mismatch:} optimizing accuracy on imbalanced data instead of recall/$F_1$/PR-AUC.
  \item \textbf{Unscaled features:} breaks distance and gradient-based methods.
  \item \textbf{Wrong validation protocol:} random splits on time series/grouped data when temporal/group leakage exists.
  \item \textbf{Overtuning CV:} too many repeated tweaks on same validation set causes validation overfit.
\end{itemize}

\subsection{Quick Recipes}
\begin{itemize}
  \item \textbf{Small tabular dataset:} start with logistic/linear regression, trees, random forest, gradient boosting.
  \item \textbf{Distance meaningful + little training time:} try $k$-NN after scaling.
  \item \textbf{High-dimensional sparse text:} linear models and linear SVMs are very strong baselines.
  \item \textbf{Images with modest data:} transfer learning with CNNs usually beats handcrafted pipelines.
  \item \textbf{Need interpretability:} linear models and shallow trees first.
  \item \textbf{Need calibrated probabilities:} logistic regression or post-hoc calibration.
\end{itemize}

\section{Rapid Topic Review}
\subsection{Introduction to ML}
\begin{itemize}
  \item ML learns from experience $E$ on tasks $T$ under performance measure $P$.
  \item Core inversion from classical programming: data + targets produce the program/model.
  \item Always separate \textbf{representation choice} from \textbf{model choice}.
\end{itemize}

\subsection{ML Tasks}
\begin{itemize}
  \item Classification predicts a discrete region/label; regression predicts a continuous response.
  \item Unsupervised learning uses similarity structure instead of labels.
  \item RL handles sequential decision-making with delayed rewards and exploration.
\end{itemize}

\subsection{Supervised Workflow}
\begin{itemize}
  \item Inspect class balance, missingness, scale, outliers, and leakage sources before modeling.
  \item Hypothesis classes define which patterns are even representable.
  \item Consistency on train is not enough; validation determines whether the design worked.
\end{itemize}

\subsection{Generalization Error}
\begin{itemize}
  \item Empirical risk is computable; true risk is an expectation under the unknown data distribution.
  \item Validation/test loss are statistical estimates of generalization error.
  \item Regularization, model-capacity control, and better data help shrink the train--test gap.
\end{itemize}

\subsection{Regression and Cross-Validation}
\begin{itemize}
  \item End-to-end regression recipe: choose features, choose loss, fit parameters, validate complexity.
  \item Overfitting often appears after a sweet spot in model complexity or training duration.
  \item $K$-fold CV reduces split variance compared to a single hold-out split.
\end{itemize}

\subsection{Ordinary Least Squares}
\begin{itemize}
  \item Design matrix rows are examples, columns are features plus an optional bias column.
  \item The normal equations solve least squares exactly when $X^\top X$ is invertible.
  \item OLS is sensitive to outliers and multicollinearity.
\end{itemize}

\subsection{$k$-NN}
\begin{itemize}
  \item Instance-based/lazy learner: defer computation until query time.
  \item Local method: prediction quality depends on the neighborhood geometry.
  \item Works best when a meaningful distance metric exists and data dimension is moderate.
\end{itemize}

\subsection{Outliers}
\begin{itemize}
  \item Outliers can be global, contextual, or collective.
  \item Visual tools: histograms, box plots, scatter plots.
  \item Robust methods prefer medians/MAD/IQR over means/standard deviations.
\end{itemize}

\subsection{Scaling}
\begin{itemize}
  \item Standardization centers at 0 with unit variance; robust scaling centers at median with IQR scale.
  \item Fit scalers on train only, then transform validation/test with the same parameters.
  \item Scaling changes optimization geometry and distance meaning, not the raw information content.
\end{itemize}

\subsection{Correlation and Polynomial Features}
\begin{itemize}
  \item One-hot is the standard encoding for nominal categories.
  \item Polynomial features convert nonlinear input relationships into a linear-in-weights problem.
  \item Correlation matrices are diagnostic tools, not causal explanations.
\end{itemize}

\subsection{Image Data I: Representation and Filters}
\begin{itemize}
  \item Images can be summarized globally (means, histograms) or locally (patches, filters).
  \item Filtering computes weighted local sums; kernels detect blur, sharpening, and edges.
  \item Gradient features capture boundaries and texture changes.
\end{itemize}

\subsection{Image Data II: Thresholds and Local Features}
\begin{itemize}
  \item Thresholding turns grayscale into binary masks for segmentation.
  \item Harris finds corners where intensity changes strongly in two directions.
  \item SIFT and ORB are classic local descriptors for matching and geometric correspondence.
\end{itemize}

\subsection{Image Data III: HOG and Classical Pipelines}
\begin{itemize}
  \item HOG summarizes local edge directions with normalization for illumination robustness.
  \item Template matching is simple but brittle under scale/rotation changes.
  \item Classical vision pipeline: preprocess, extract features, normalize, classify with a standard ML model.
\end{itemize}

\subsection{Linear Models}
\begin{itemize}
  \item Linear models produce a score $\w^\top\phi(\x)+b$ on features or transformed features.
  \item Regression/classification differ mostly in output interpretation and loss.
  \item ERM + regularization is the common training principle.
\end{itemize}

\subsection{Gradient Methods}
\begin{itemize}
  \item Gradients point uphill; gradient descent steps downhill.
  \item Batch size trades variance against compute efficiency.
  \item Adaptive optimizers are practical defaults, but learning-rate choice still matters.
\end{itemize}

\subsection{SVMs}
\begin{itemize}
  \item SVMs maximize margin, not probability fit.
  \item Soft-margin parameter $C$ balances margin width against training violations.
  \item Kernels implicitly lift data into richer feature spaces.
\end{itemize}

\subsection{Logistic Regression}
\begin{itemize}
  \item Logistic regression is a linear score model with probabilistic interpretation.
  \item Cross-entropy is the natural loss from Bernoulli/multinomial likelihoods.
  \item Decision thresholds can be tuned independently of model training.
\end{itemize}

\subsection{Decision Trees and Ensembles}
\begin{itemize}
  \item Trees greedily choose splits maximizing impurity/variance reduction.
  \item Deep trees memorize; pruning and ensembling reduce variance.
  \item Random forests usually outperform a single tree while sacrificing some interpretability.
\end{itemize}

\subsection{MLP Basics}
\begin{itemize}
  \item Hidden layers act as learned feature extractors.
  \item Nonlinear activations are what make depth useful.
  \item Width/depth choices control expressivity and optimization behavior.
\end{itemize}

\subsection{MLP Training and Softmax}
\begin{itemize}
  \item Backpropagation efficiently computes all gradients via repeated chain rule application.
  \item Softmax + cross-entropy is the standard multi-class pairing.
  \item Training curves should be read jointly: train loss, validation loss, and validation metric.
\end{itemize}

\subsection{CNNs}
\begin{itemize}
  \item Local connectivity and weight sharing drastically cut parameter count.
  \item Pooling/downsampling trades spatial precision for abstraction and efficiency.
  \item Transfer learning is often the strongest baseline on moderate image datasets.
\end{itemize}

\subsection{Transformers and GenAI}
\begin{itemize}
  \item Self-attention lets each token condition on the whole context.
  \item BERT-style encoders are for understanding; GPT-style decoders are for generation.
  \item Prompting changes conditional generation without changing model weights.
\end{itemize}

\subsection{Practical LLM/GAN Notes}
\begin{itemize}
  \item Tokenization quality strongly affects downstream NLP behavior.
  \item Fine-tuning adapts pretrained representations instead of learning from scratch.
  \item GAN training is unstable because generator and discriminator co-evolve.
\end{itemize}

\section{Exam-Style Facts to Remember}
\begin{enumerate}
  \item Simpler models often win with limited data; flexibility needs regularization and validation.
  \item Correlation measures linear association only; nonlinear dependence can hide behind $r\approx 0$.
  \item Tree-based models are usually invariant to monotone rescaling; distance/margin/gradient models are not.
  \item One-hot encoding is safe for nominal categories; label encoding is not unless order is real.
  \item Ridge shrinks all coefficients; lasso can perform feature selection.
  \item Logistic regression is linear in features but nonlinear in probability via the sigmoid.
  \item SVM support vectors are the informative training points; non-support vectors do not affect the margin solution much.
  \item Deep nets learn features automatically, but good optimization and regularization are essential.
  \item CNNs beat dense nets on images because locality + weight sharing encode the right inductive bias.
  \item Transformers succeed because attention models long-range dependence in parallel.
\end{enumerate}

\section{Formula Bank}
\[
L_{\mathrm{emp}}=\frac{1}{N}\sum_{i=1}^N \ell\big(f(\x_i),y_i\big)
\]
\[
\mathrm{SSE}=\sum_i (y_i-\hat y_i)^2,
\qquad
\mathrm{MSE}=\frac{1}{N}\sum_i (y_i-\hat y_i)^2,
\qquad
\mathrm{RMSE}=\sqrt{\mathrm{MSE}}
\]
\[
\mathrm{MAE}=\frac{1}{N}\sum_i |y_i-\hat y_i|,
\qquad
z=\frac{x-\mu}{\sigma},
\qquad
x'_{\mathrm{std}}=\frac{x-\mu}{\sigma}
\]
\[
\sigma(z)=\frac{1}{1+e^{-z}},
\qquad
\log\frac{p}{1-p}=\w^\top\x+b
\]
\[
\mathrm{softmax}(z_i)=\frac{e^{z_i}}{\sum_j e^{z_j}},
\qquad
\ell_{\mathrm{CE}}=-\sum_k y_k\log p_k
\]
\[
\ell_{\mathrm{hinge}}=\max(0,1-yf(\x)),
\qquad
\gamma_i=\frac{y_i(\w^\top\x_i+b)}{\|\w\|}
\]
\[
\hat{\w}_{\mathrm{OLS}}=(X^\top X)^{-1}X^\top\y,
\qquad
\w_{t+1}=\w_t-\eta\nabla L(\w_t)
\]
\[
\hat y = w_0+\sum_j w_j x_j,
\qquad
r_{xy}=\frac{\mathrm{cov}(x,y)}{\sigma_x\sigma_y}
\]
\[
H(S)=-\sum_c p_c\log_2 p_c,
\qquad
G(S)=1-\sum_c p_c^2
\]
\[
IG(S,\text{split}) = H(S)-\sum_j \frac{|S_j|}{|S|}H(S_j)
\]
\[
\text{Accuracy}=\frac{TP+TN}{TP+TN+FP+FN},
\quad
\text{Precision}=\frac{TP}{TP+FP},
\quad
\text{Recall}=\frac{TP}{TP+FN}
\]
\[
F_1=\frac{2PR}{P+R},
\qquad
\text{IQR}=Q_3-Q_1,
\qquad
x'_{\mathrm{robust}}=\frac{x-\mathrm{median}}{Q_3-Q_1}
\]
\[
L_{\mathrm{ridge}}=\|X\w-\y\|_2^2+\lambda\|\w\|_2^2,
\qquad
L_{\mathrm{lasso}}=\|X\w-\y\|_2^2+\lambda\|\w\|_1
\]
\[
\delta^{(\ell)} = \big(W^{(\ell+1)\top}\delta^{(\ell+1)}\big)\odot \phi'(a^{(\ell)})
\]
\[
\frac{\partial L}{\partial W^{(\ell)}} = \delta^{(\ell)} h^{(\ell-1)\top}
\]
\[
\mathrm{Attention}(Q,K,V)=\mathrm{softmax}\!\left(\frac{QK^\top}{\sqrt{d_k}}\right)V
\]
\[
\text{conv output size} = \left\lfloor\frac{n+2p-f}{s}\right\rfloor+1
\]

\end{document}
